%% THERMODYNAMIQUE %%
%
\section{Composantes du vent \`a partir des vitesses "avion" dans le rep\`ere terrestre}

\itemstyle{ Op\'erateur~: } \operfmt{CalcVent\_repgeo}
\index{CalcVent\_repgeo}

%%%%%OBJET%%%%%
\itemstyle{ Objet~: }
Cet op\'erateur calcule les 3 composantes du vecteur vent. La vitesse
de l'air par rapport \`a l'avion est transform\'ee vers le
rep\`ere g\'eographique terrestre 
(ENU : axes orient\'es vers l'Est, le Nord, et verticalement).
Puis, le vent est calcul\'e dans ce rep\`ere par combinaison avec la
vitesse de l'avion par rapport au sol.
Cet op\'erateur tient compte de la distance entre le point
de mesure de la vitesse de l'air par rapport \`a l'avion
et le point de mesure de la vitesse avion par rapport au sol.
Mais le calcul suppose que ces deux points sont situ\'es
sur l'axe longitudinal de la plate-forme mobile.

%%%%%ARGUMENTS%%%%
\bigskip
%%%%%ENTREES%%%%
\begin{tabarg}
\multicolumn{3}{l}{\titretab{ Entr\'ees}}\\
\hline
$L$ & Coeff., float& Distance alg\'ebrique entre le point de mesure de la vitesse 
de l'air par rapport \`a l'avion, et le point de mesure de la 
vitesse avion par rapport au sol ($m$)\\
$V_p$ & Vecteur, {\tt float} & Vitesse propre ($m/s$) \\
$\alpha$ & Vecteur, {\tt float} & Angle d'attaque, positif lorsque le vent relatif vient de sous l'avion ($radians$) \\
$\beta$ & Vecteur, {\tt float} & Angle de d\'erapage, positif lorsque le 
vent relatif vient de tribord ($radians$)\\
$V_e$ & Vecteur, {\tt float} & Composante zonale de la vitesse de l'avion par rapport 
au sol ($m/s$)\\
$V_n$ & Vecteur, {\tt float} & Composante m\'eridienne de la vitesse
de l'avion par rapport au sol ($m/s$)\\ 
$V_v$ & Vecteur, {\tt float} & Composante verticale de la vitesse de 
l'avion par rapport au sol ($m/s$)\\
$Vr_y$ & Vecteur, {\tt float} & Vitesse de rotation suivant l'axe Y, positif quand $\theta$ cro\^{\i}t($rad/s$) \\
$Vr_z$ & Vecteur, {\tt float} & Vitesse de rotation suivant l'axe Z, positif quand $\psi$ cro\^{\i}t ($rad/s$) \\
$\theta$ & Vecteur, {\tt float} & Angle de tangage de l'avion, positif quand
l'avion est cabr\'e ($radians$)\\
$\phi$ & Vecteur, {\tt float} & Angle de roulis de l'avion, positif quand
l'aile tribord est enfonc\'ee ($radians$)\\
$\psi$ & Vecteur, {\tt float} & Cap g\'eographique de l'avion ($radians$) \\
\end{tabarg}

\medskip
%%%%%SORTIES%%%%
\begin{tabarg}
\multicolumn{3}{l}{\titretab{ Sorties}}\\
\hline
$U$ & Vecteur, {\tt float} & Composante zonale du vent (Ouest en Est) ($m/s$)\\
$V$ & Vecteur, {\tt float} & Composante m\'eridienne du vent (Sud au Nord) ($m/s$)\\
$W$ & Vecteur, {\tt float} & Composante verticale du vent (Bas en Haut) ($m/s$)\\
\end{tabarg}



%%%%%FORMULE%%%%
\itemstyle{ Formule }

Les composantes de la vitesse  de l'air
par rapport � l'avion, corrig�es de la vitesse d'entrainement
de rotation, sont d'abord calcul�es dans le rep�re avion~: 
\begin{eqnarray*}
U_a & = & - V_p.D^{-1} \\
V_a & = & - V_p.D^{-1}.\tan\beta + L.Vr_z\\
W_a & = & - V_p.D^{-1}.\tan\alpha - L.Vr_y
\end{eqnarray*}

avec~: $ D = \sqrt{1+\tan^2\alpha+\tan^2\beta}$

Une rotation fournit les componsantes 
dans le rep�re g�ographique :
\newcommand\relphantom[1]{\mathrel{\phantom{#1}}}
\begin{equation*}\begin{split}
U_{ag} &= U_a.\sin\psi.\cos\theta +  
          V_a.(\cos\psi.\cos\phi + \sin\psi.\sin\theta.\sin\phi) \\
      &\relphantom{=} +  W_a.(\sin\psi.\sin\theta.\cos\phi - \cos\psi.\sin\phi) \\   
V_{ag} &= U_a.\cos\psi.\cos\theta - 
          V_a.(\sin\psi.\cos\phi - \cos\psi.\sin\theta.\sin\phi) \\
      &\relphantom{=}  + W_a.(\cos\psi.\sin\theta.\cos\phi + \sin\psi.\sin\phi) \\
W_{ag} &= U_a.\sin\theta - 
          V_a.\cos\theta.\sin\phi -
          W_a.\cos\theta.\cos\phi
\end{split}\end{equation*}
%\begin{equation*}
%\begin{pmatrix}
%U_{ag} \\ V_{ag} \\ W_{ag} \\
%\end{pmatrix} =
%\begin{pmatrix}
%\sin\psi.\cos\theta & \cos\psi.\cos\phi + \sin\psi.\sin\theta.\sin\phi & 0 \\
%\cos\psi.\cos\theta & \sin\psi.\cos\phi - \cos\psi.\sin\theta.\sin\phi & 0 \\
%\sin\theta & 1 & 0 \\
%\end{pmatrix}
%\begin{pmatrix}
%U_{a} \\ V_{a} \\ W_{a} \\
%\end{pmatrix}
%\end{equation*}

Enfin, le vent est la somme des vitesses
de l'air par rapport � l'avion et de l'avion par rapport
au sol~:
\begin{equation*}
\begin{pmatrix}U \\ V \\ W  \end{pmatrix} = 
\begin{pmatrix}U_{ag} \\ V_{ag} \\  W_{ag} \end{pmatrix} + 
\begin{pmatrix} V_e \\ V_n \\ V_v  \end{pmatrix}
\end{equation*}


%%%%%AUTEUR%%%%%
\itemstyle{ Auteur~:}
CNRM/GMEI/TRAMM

%%%%%REFERENCES%%%%%
\itemstyle{ R\'ef\'erences~:}
Bulletin NCAR-RAF Nr 23 \cite{RAF-23}.
